%==============================================================================
%  Joint Transport--Relay Co-Design for UAV Emergency Logistics in Mountainous
%  Terrain:  Handover-Aware Feasibility, Channel-Model Sensitivity, and the
%  Price of Partitioning
%
%  English SCI manuscript (draft).  All numbers are macros generated by
%  code/make_numbers_en.py from the result files; see numbers_en_sources.txt
%  for the provenance of every macro.
%==============================================================================
\documentclass[11pt]{article}

\usepackage[margin=2.4cm]{geometry}
\usepackage{amsmath,amssymb,amsthm}
\usepackage{booktabs}
\usepackage{graphicx}
\usepackage{siunitx}
\usepackage{enumitem}
\usepackage[hidelinks]{hyperref}
\usepackage{xcolor}
\usepackage{caption}
\usepackage{algorithm}
\usepackage{algpseudocode}
\algrenewcommand\algorithmicrequire{\textbf{Input:}}
\algrenewcommand\algorithmicensure{\textbf{Output:}}
\floatname{algorithm}{Algorithm}
\captionsetup{font=small}

% ---- \texttt 断行：正文引用 code/make_numbers_en.py 这类长路径，
% IEEEtran 双栏窄栏下无法断行会 overfull。改为每个字符后允许换行。
% 必须先用 \makeatletter 定义辅助宏、再 \renewcommand，否则展开时未定义。
\makeatletter
\def\dsh@tt#1{\dsh@ttloop#1\dsh@end}
\def\dsh@ttloop#1{\ifx#1\dsh@end\else\dsh@ttemit{#1}\expandafter\dsh@ttloop\fi}
\def\dsh@ttemit#1{#1\penalty\z@\hskip\z@ plus 0.02em\relax}
\renewcommand{\texttt}[1]{{\ttfamily\hyphenchar\font=-1 \dsh@tt{#1}}}
\makeatother

% ---- 表格自适应列宽 ---------------------------------------------------------
% 同一个 body 要能在三种版式下编译：单栏 article / elsarticle preprint，
% 以及 IEEEtran 的双栏窄栏。表格按单栏宽度排就会在双栏里溢出（右列被裁掉），
% 所以所有表格统一走 \fitwide：把表格缩放到当前版心的宽度。
%   * 单栏版式：表格本身通常窄于版心，缩放系数 > 1 会被 \resizebox 忽略（只缩不放时
%     用 \fitwide 的 \ifdim 分支），因此不会把小表拉大；
%   * 双栏版式：在 table* 里 \linewidth = \textwidth（跨双栏），表格得到足够宽度。
\newsavebox{\wearfitbox}
\newcommand{\fitwide}[1]{%
  \sbox{\wearfitbox}{#1}%
  \ifdim\wd\wearfitbox>\linewidth
    \resizebox{\linewidth}{!}{\usebox{\wearfitbox}}%
  \else
    \usebox{\wearfitbox}%
  \fi
}



% ---- 结果文件缺失时给出显式提示，而不是静默留空 ----------------------------
% 依序尝试：调用时给出的路径 → 当前目录 → paper_en/ → 上一级的 paper_en/ → 上一级。
% 这样在三种布局下都能编：工作区（在 D题/ 下编 paper_en/main.tex）、
% paper_en/ 内直接编译、以及交付包（06_英文SCI稿/ 内编译）。
\newcommand{\InputRes}[1]{%
  \IfFileExists{#1}{\input{#1}}{%
    \IfFileExists{paper_en/#1}{\input{paper_en/#1}}{%
      \IfFileExists{../paper_en/#1}{\input{../paper_en/#1}}{%
        \IfFileExists{../#1}{\input{../#1}}{%
          \typeout{== WARNING: resource #1 NOT FOUND ==}%
          \par\noindent\fbox{\parbox{0.9\linewidth}{\centering Missing resource
          \texttt{#1}; run \texttt{python code/make\_numbers\_en.py}}%
          }\par
        }%
      }%
    }%
  }%
}
% ---- 插图：优先用矢量 PDF，回退到 300 dpi PNG -------------------------------
% \Dfig[宽度]{不带扩展名的图名}{图注}{label}
% 为什么带扩展名解析：数据图由 make_figs_en.py / make_fig_theory.py 以
% **矢量 PDF + 300 dpi PNG** 双份导出。双栏排版把图缩到约 8.8 cm 宽，
% 只有矢量版本在任意缩放下锐利；而 DEM 晕渲底图本身是栅格，转矢量只会
% 撑大 PDF 不增清晰度，因此**只有 PNG**，这里靠回退逻辑自动取用。
\newcommand{\Dfigfile}[3]{%
  \IfFileExists{#3#2.pdf}{\includegraphics[width=#1\linewidth]{#3#2.pdf}}{\IfFileExists{#3#2.png}{\includegraphics[width=#1\linewidth]{#3#2.png}}{\IfFileExists{#3#2.eps}{\includegraphics[width=#1\linewidth]{#3#2.eps}}{\fbox{\parbox{0.7\linewidth}{\centering Figure \texttt{#2} not yet generated; run \texttt{python code/make\_figs\_en.py}}}}}}%
}
% 插图入口：\Dfig[宽度]{图名（不带扩展名）}{图注}{label}
% 图名一律省略扩展名——由 \Dfigfile 按 pdf → png → eps 的顺序挑，
% 这样"数据图用矢量、DEM 晕渲用 300 dpi 栅格"这件事由文件是否在场决定，
% 正文不需要知道哪张是矢量。
\newcommand{\Dfig}[4][0.92]{%
  \begin{figure}[htbp]\centering
    \Dfigfile{#1}{#2}{figs/}\caption{#3}\label{#4}%
  \end{figure}%
}

% ---- 标题统一：同一字体族、不用全大写 ----------------------------------------
% 三种版式对 \section 的默认处理不一样：
%   * article（中立版）：\Large\bfseries，本来就保留大小写；
%   * IEEEtran：\normalfont\normalsize\centering\scshape —— 小体大写，
%     于是 "System model" 印成 "SYSTEM MODEL"；
%   * elsarticle：\large\bfseries\scshape 或默认 article 样式。
% 结果是"章标题全大写、节标题正常大小写"，两套样式并存；审稿时看着像格式不统一。
% 这里统一为：同一字体族（正文罗马体）+ 同一字重（粗体）+ 保留**句子大小写**，
% 三级标题（section / subsection / subsubsection）保持一致，不再出现全大写。
% 位置参数沿用各类的默认值，避免破坏正文与标题的间距。
\makeatletter
\providecommand{\dsh@section@spacing}{%
  {1\baselineskip plus 0.25\baselineskip minus 0.25\baselineskip}%
  {0.7ex plus 1ex minus 0ex}}
\@ifclassloaded{IEEEtran}{%
  \def\section{\@startsection{section}{1}{\z@}%
    {1.0\baselineskip plus 0.25\baselineskip minus 0.25\baselineskip}%
    {0.4\baselineskip}%
    {\normalfont\normalsize\centering\bfseries}}%
  \def\subsection{\@startsection{subsection}{2}{\z@}%
    {1.0\baselineskip plus 0.25\baselineskip minus 0.25\baselineskip}%
    {0.3\baselineskip}%
    {\normalfont\normalsize\itshape}}%
  \def\subsubsection{\@startsection{subsubsection}{3}{\z@}%
    {0.8\baselineskip plus 0.2\baselineskip minus 0.2\baselineskip}%
    {0.2\baselineskip}%
    {\normalfont\normalsize\itshape}}%
}{%
  % article（中立版）与 elsarticle：只需去掉 small-caps，字族字重沿用默认
  \@ifundefined{elsarticle@type}{%
    \renewcommand{\section}{\@startsection{section}{1}{\z@}%
      {-3.5ex plus -1ex minus -0.2ex}{2.3ex plus 0.2ex}%
      {\normalfont\Large\bfseries}}%
  }{%
    \renewcommand{\section}{\@startsection{section}{1}{\z@}%
      {-3.5ex plus -1ex minus -0.2ex}{2.3ex plus 0.2ex}%
      {\normalfont\large\bfseries}}%
  }%
}
\makeatother

% ---- 公式自适应字号 ---------------------------------------------------------
% 有几个长行间公式在单栏版式下正好放得下，但在 IEEEtran 的单栏里会溢出
% （\hbox too wide，压到栏外）。\eqfit 用 graphicx 的 \width 直接判断自然宽度：
% 放不下才按比例缩到版心，放得下就维持 100% 原样，因此单栏版式排版不变。
% 注意整条公式只有一个 \resizebox，不要在 \ifdim 分支里各写一份——那样宏内的
% 换行会让 \setbox 少吃掉一个空格并把行号推后一行。
%
% 另外两处细节是排版事故的来源，必须保留：
%   * 开头的 \par 把前一段**收尾**。否则当 \eqfit 紧跟在 \begin{proof} 之后时，
%     "Proof." 之后那一行没有可断处，TeX 会把极少几个词撑满整栏——看起来就是
%     "Proof.        The        cruise        altitude        is" 那种诡异间距。
%   * 结尾的 \par 把 center 段落封口，避免后文与公式之间产生额外空行。
\newcommand{\eqfit}[1]{%
  \par\nobreak
  \begin{center}%
  \resizebox{\ifdim\width>\linewidth\linewidth\else\width\fi}{!}{$\displaystyle #1$}%
  \end{center}%
  \par\nobreak
}

% ---- 浮动体间距：避免图与正文小标题贴得太近（双栏下尤其明显）----------------
\setlength{\textfloatsep}{12pt plus 3pt minus 2pt}
\setlength{\floatsep}{12pt plus 3pt minus 2pt}
\setlength{\intextsep}{12pt plus 3pt minus 2pt}
\renewcommand{\topfraction}{0.92}
\renewcommand{\bottomfraction}{0.85}
\renewcommand{\textfraction}{0.06}
\renewcommand{\floatpagefraction}{0.82}

% ---- 定理类环境：统一为标准版式，标题独立成行 --------------------------------
% amsthm 默认把**标题与正文放在同一段**，于是标题也参与该行的两端对齐。当正文
% 以一条长的不可断行公式开头时（例如 Corollary 5 的
% $D=\sum_k\max(0,w-\mathrm{span}(k-w))$），这一行没有可断处，TeX 只能把标题里的
% 几个空格撑开去填满整栏——双栏版式下就变成
% "Corollary        5        (Deficiency        functional)."
% 这类忽宽忽窄的样子（单栏时栏宽大、不易察觉，双栏时非常明显）。
%
% 修法：把标题与正文之间的分隔（\newtheoremstyle 的第 7 个参数）改成
% "\newline"，让**标题独占一行**（"By name" 风格）。标题行只含标题本身，
% 没有可拉伸的余地；正文另起一行，两端对齐不再受公式长度牵制。
% 标题之后不再加句点，这也与常见的 "By name" 定理样式一致。
\newtheoremstyle{dshplain}%
  {\topsep}{\topsep}%
  {\itshape}%   正文
  {}%           缩进
  {\bfseries}%  标题
  {}%           标题后标点：标题独立成行，不加句点
  {\newline}%   标题与正文之间：换行
  {}%           标题格式
\newtheoremstyle{dshdefn}%
  {\topsep}{\topsep}%
  {\normalfont}% 定义/假设用直立体，与定理类区分
  {}%
  {\bfseries}%
  {}%
  {\newline}%
  {}
\theoremstyle{dshplain}
\newtheorem{theorem}{Theorem}
\newtheorem{proposition}[theorem]{Proposition}
\newtheorem{corollary}[theorem]{Corollary}
\theoremstyle{dshdefn}
\newtheorem{assumption}[theorem]{Assumption}
\newtheorem{definition}[theorem]{Definition}

\newcommand{\tbd}{\texttt{TBD}}
\newcommand{\y}[1]{\textsc{#1}}

% ---- 紧急拉伸：12pt 单栏版式下有成串的行内数学（FSPL 不等式、
% Fresnel-Kirchhoff 参数等），这些数学整体不可断字，偶尔会把一行撑出版心。
% \emergencystretch 只在原本会 overfull 的行上起作用，不影响正常排版。
% 放在这里是为了同时进入三种版式（期刊变体只取"共享段"）。
\setlength{\emergencystretch}{3em}

\InputRes{numbers_en.tex}

\title{Supplementary Material}
\date{\today}

\begin{document}
\maketitle

\section*{Supplementary material}\label{supp}

This supplement collects the proofs and the pseudocode listings that are
omitted from the main text for length. Equation and theorem numbering
follows the main text; every item below is flagged at its point of use
there.

\subsection*{Proof of \texttt{cor:continuous}}
\begin{proof}
Identical to Theorem~\ref{thm:handover}; discreteness is used nowhere in that
argument.
\end{proof}

\subsection*{Proof of \texttt{prop:demand}}
\begin{proof}
A point that covers $D'$ on an interval covers $D$ on it, so every span value
weakly decreases; the deficiency is a sum of pointwise non-increasing terms.
Any schedule that is valid for $D'$ is valid for $D$, since validity is a
covering condition. The energy statement is immediate from
Assumption~\ref{asm:energy}.
\end{proof}

\subsection*{Proof of \texttt{prop:energy}}
\begin{proof}
Let $A_i\subseteq[t_0,t_{n-1}]$ be the set of instants at which relay $i$ is
airborne and therefore able to provide coverage. Continuous coverage means
$[t_0,t_{n-1}]=\bigcup_{i=1}^{N}A_i$, so $T\le\sum_i|A_i|$. By
Assumption~\ref{asm:energy} each maximal interval of $A_i$ has length at most
$T^{\mathrm{cap}}$, but $A_i$ may consist of $m_i>1$ such intervals, so
$|A_i|\le m_iT^{\mathrm{cap}}$ and the decomposition alone yields only
$T\le T^{\mathrm{cap}}\sum_i m_i$, which says nothing about $N$: a single relay
recharging repeatedly is not excluded by Assumption~\ref{asm:energy}. The bound
therefore requires one further hypothesis, which we state explicitly and verify on
our instance rather than leave implicit:

\smallskip
\noindent\emph{(H) the coverage intervals of distinct relays can be taken
pairwise disjoint.}
\smallskip

\noindent Under (H) at most one relay is covering at any instant, the airtime of
distinct relays does not overlap in the covering role, and
$T=\sum_i|A_i|\le\sum_iT^{\mathrm{cap}}=N\,T^{\mathrm{cap}}$, giving
$N\ge\lceil T/T^{\mathrm{cap}}\rceil$.

\smallskip
\noindent On our instance (H) is exactly what the mission demands rather than a
convenience: coverage must be sustained at $\ENqCells$ demand cells spanning
$T=\ENTs$\,s, and the verified schedule at $N=\ENrbGreedy$ uses
$\ENgrNThreeMOneSorties$ sorties with SOC minimum $\ENgrNThreeMOneSoc\%$, i.e.\ the
relays \emph{hand the covering role over} instead of overlapping it. We therefore
report $N\ge\ENboundEnergy$ as a bound \emph{conditional on (H)}, and we do not
rely on it for any fleet-size claim that the constructive schedule does not
already support.
\end{proof}

\subsection*{Proof of \texttt{prop:lb}}
\begin{proof}
Fix a feasible schedule and a phase $\Phi_i$. Let $P_i$ be the set of hover
points \emph{actually occupied} during $\Phi_i$ by the $N$ relays. Condition (V)
requires that the union of the points occupied at each instant covers that
instant's demand, and each occupied point is one of the at most $N$ relays. Hence
$P_i$ is a set of at most $N$ points covering all demand in $\Phi_i$, so
$|P_i|\ge c_i$ by definition of $c_i$. Since $|P_i|\le N$, we get $N\ge c_i$, and
$i$ was arbitrary. (Note that this bound does not assume the relays are static
across the phase --- only that at most $N$ points are occupied at any instant, so
a relay that changes point inside a phase only enlarges $P_i$.)
\end{proof}

\subsection*{Proof of \texttt{prop:lipschitz}}
\begin{proof}
The cruise altitude is $\max\{\text{line maximum},\ \text{endpoint altitudes}\}+50$,
a maximum of $1$-Lipschitz and $M$-Lipschitz functions of the endpoints, hence
$(1+M)$-Lipschitz; the climb and descent heights are this quantity minus a
constant, and the horizontal distance is $1$-Lipschitz by the triangle
inequality. Composing with the linear time map gives the constant $L$.
\end{proof}

\subsection*{Proof of \texttt{prop:monotone}}
\begin{proof}
Keep the schedule of the $N$ relays and let the additional relay remain at the
depot; it neither helps nor hinders.
\end{proof}

\subsection*{Proof of \texttt{prop:phase}}
\begin{proof}
{\sloppy\emph{Forward.} Given a phase-stationary schedule, assign to each relay
and each phase the point at which that relay hovers during that phase.}\par A relay serves at
most one point per phase by hypothesis, and every incidence is covered because the
schedule covers every instant. It remains to check (ii). If a relay uses different
points in two consecutive phases, it must relocate; by
Assumption~\ref{asm:handover} the relocation occupies an interval of length at
least $W(a,b)$ during which the relay covers nothing, and by
Assumption~\ref{asm:serial} that interval may not overlap another relay's
relocation. Since the relay must be in place for the whole of phase $k+1$, the
interval must finish by the start of phase $k+1$; since it must have left its
phase-$k$ point, it starts no earlier than the end of phase $k$ at which it last
held that point. Hence $W(a,b)$ is at most the length of the inter-phase gap, which
is condition (ii). Hence the relaxed model admits the induced assignment.
{\sloppy\emph{Converse.} The relaxation drops the requirement that the
\emph{stationary} relay hold a single point throughout the moving relay's
blackout window; on our instance this extra requirement is what makes $N=2$
infeasible while the relaxation is satisfiable.}\par
\end{proof}

\subsection*{Proof of \texttt{prop:structural}}
\begin{proof}
Assumption~\ref{asm:serial} makes the blackout windows of all relocations
pairwise disjoint; each has length at least $w_{\min}$ and each completion lies in
$[0,T]$, so the completions are separated by at least $w_{\min}$ and
$r\,w_{\min}\le T$. It remains to connect $N$ to $r$. By the premise, no single
configuration (one hover point per relay) covers the whole profile, so for each
relay $i$ there is a demand cell that $i$'s point fails to cover; covering that
cell therefore requires some relay to be elsewhere, i.e.\ to relocate. This holds
for every relay, so all $N$ relocate and $r\ge N$. Combining gives
$N\le r\le 1+\lfloor T/w_{\min}\rfloor$.
\end{proof}

\subsection*{Proof of \texttt{thm:one}}
\begin{proof}
A single relay cannot relocate without leaving the demand uncovered, because
relocation always costs a strictly positive blackout. That positivity is part
(ii) for M2 ($W(a,a)=t_{\mathrm{link}}>0$, so even a no-op relocation blacks the
link out) and the closing sentence of Proposition~\ref{prop:W} for M1
($W(a,a)>0$ there as well); note that it is \emph{not} part (iii), which is the
triangle bound and says nothing about sign. Hence its hover
point must be constant throughout $[t_{\mathrm{f}},t_{\ell}]$, which is exactly
condition 1, and the sortie duration is exactly the expression in condition 2.
\end{proof}

\subsection*{Pseudocode for \texttt{sec:feedback}}
\begin{algorithm}[htbp]
\caption{Co-design feedback loop (threshold accepting on departure instants)}
\label{alg:code}
\begin{algorithmic}[1]
\Require frozen batching, routes, aircraft types and assignments; feasible
departure window $[l_i,u_i]$ and duration $d_i$ for each sortie $i$ (so the
latest departure that still finishes inside the window is $u_i-d_i$); deficiency
oracle $D(\cdot)$; step scale $\sigma$
\Ensure best departure vector and its deficiency
\State $s\gets s^{0}$;\ \ $D^{*}\gets D(s^{0})$;\ \ $s^{*}\gets s^{0}$
\For{$r=1$ to $R$} \Comment{restarts}
  \State $s\gets$ perturb($s^{0}$)
  \For{$\mathrm{it}=1$ to $I$}
    \State pick $i$; \ $s'\gets s$ with $s_i\gets\mathrm{clip}(s_i+\mathcal{N}(0,\sigma),\,l_i,\,u_i-d_i)$
    \State $T\gets T_0\,(1-\mathrm{it}/I)$
    \If{$D(s')\le D(s)$ \textbf{or} $\mathrm{rand}()<\exp\!\bigl(-(D(s')-D(s))/T\bigr)$}
      \State $s\gets s'$
    \EndIf
    \If{$D(s)<D^{*}$} \State $D^{*}\gets D(s)$;\ \ $s^{*}\gets s$ \EndIf
  \EndFor
\EndFor
\State \Return $s^{*},D^{*}$
\end{algorithmic}
\end{algorithm}

\subsection*{Pseudocode for \texttt{thm:exact}}
\begin{algorithm}[htbp]
\caption{Handover-aware relay feasibility (stated without pruning; the
runnable implementation adds the pruning reported in
Section~\ref{sec:reproducible})}
\label{alg:dp}
\begin{algorithmic}[1]
\Require demand cells $t_0<\dots<t_{n-1}$ with sets $\mathcal{D}_k$; candidate
points $\mathcal{C}$ with coverage $\mathrm{cov}[k,c]$; blackout $W$; cap
$T^{\mathrm{cap}}$; fleet size $N$
\Ensure a feasible schedule, or the first cell with no feasible state
\State $S_0 \gets \bigl\{\,((\mathrm{IDLE},t_0,t_0))_{i=1}^{N}\,\bigr\}$
\For{$k=0$ to $n-1$}
  \State $S_k \gets \{\,s\in S_k:\ \bigcup_{i:\,p_i\neq\mathrm{IDLE}}\mathrm{cov}[\cdot,p_i]\supseteq\mathcal{D}_k\,\}$ \Comment{(V)}
  \If{$S_k=\emptyset$} \State \Return ``no feasible state at cell $k$''
  \EndIf
  \ForAll{$s\in S_k$, $i\in\{1..N\}$, $p'\in\mathcal{C}\cup\{\mathrm{IDLE}\}$, $p'\neq p_i$}
    \State $j \gets j_k \gets \textsc{Before}(k,\,W(p_i,p'))$ \Comment{last cell with $t_{j}\le t_k-W$; not $\lceil W/\Delta\rceil$}
    \If{(T1) $\sigma_i\le t_{j}$ \textbf{and} (T2) $\sigma_j\le t_{j}\ \forall j\neq i$}
      \If{(T3) $\bigcup_{j\neq i}\mathrm{cov}[\cdot,p_j]\supseteq\mathcal{D}_{k'}\ \forall k'\in[j+1,k-1]$}
        \If{(T4) $\bigcup_{j}\mathrm{cov}[\cdot,p_j]\supseteq\mathcal{D}_k$ \textbf{and}
             (T5) $t_k-\varrho_i^{(k)}\le T^{\mathrm{cap}}\ \forall i$}
          \State $S_{k+1}\gets S_{k+1}\cup\{\,s\ \text{with}\ (p_i,\sigma_i)\gets(p',t_k),\ \varrho_i\gets t_k\ \text{iff mode M1 or }p_i=\mathrm{IDLE}\,\}$
        \EndIf
      \EndIf
    \EndIf
  \EndFor
\EndFor
\State \Return any state path in $S_{n-1}$
\end{algorithmic}
\end{algorithm}

\end{document}
